\documentclass[UTF8]{ctexart}

\usepackage{graphicx}
\usepackage{amsmath}
\usepackage{booktabs}
\usepackage{geometry}
\usepackage{float}
\usepackage{caption}
\usepackage{times}
\usepackage{mathptmx}
\usepackage[hidelinks]{hyperref}

\geometry{a4paper, margin=2.5cm}

\title{基于DQN族深度强化学习的果园无人机喷洒路径规划方法研究}
\author{Zheng Yayu}
\date{\today}

\begin{document}

\maketitle

\begin{abstract}
面向精准农业中的果园无人机喷洒任务，路径规划算法需要在保证作业覆盖率和飞行安全的同时，尽可能缩短飞行距离、降低重复喷洒，并适应地块尺度扩大、变量喷洒需求和障碍物风险等复杂条件。针对传统规则遍历方法在复杂果园结构下路径冗余较大、难以与在线避障和喷洒决策统一优化的问题，本文基于 Aerial Autonomy Stack（AAS）原有 Gazebo 果园仿真环境，构建了一套面向无人机喷洒路径规划的 DQN 族深度强化学习实验框架。实验并未重新手工搭建果园地图，而是直接解析 AAS 中的 \texttt{apple\_orchard.sdf} 文件，自动提取目标树木、静态障碍物和农田区域，并将其转换为网格化路径规划环境。

本文将喷洒路径规划建模为马尔可夫决策过程，设计了包含无人机位置、最近未覆盖目标、覆盖进度、局部障碍状态、变量需求和动态安全值的状态表示；动作空间由上下左右离散移动构成，并在增强任务中引入喷洒开关和自动喷洒控制；奖励函数综合考虑新增覆盖、剂量满足、重复喷洒、路径长度、碰撞惩罚、动态障碍安全值 $S_v$ 和任务完成奖励。为分析不同强化学习算法的适用性，本文对 DQN、Double DQN、Dueling DQN、Rainbow DQN lite、PPO、A2C、SAC、TD3 和 DRQN 进行了分阶段对比，并设置静态树行覆盖、智能灌溉需求、动态障碍增强和整图多区块覆盖四类实验。

实验结果表明，普通 DQN 可以作为基础价值函数强基线，但并不适合作为最终完整系统的唯一主体算法。在原始树行静态覆盖任务中，Dueling DQN 表现最稳定，适合作为静态果园喷洒路径规划的强基线；在分层自动喷洒和安全控制增强任务中，本文将 DQN、DRQN、Dueling DQN 和 Rainbow DQN lite 补充到 5 个随机种子，DQN 达到 80.0\% 成功率，DRQN 的平均需求满足率最高，为 $80.9\%\pm14.7\%$，四种算法的规划碰撞和动态碰撞均为 0。消融实验进一步表明，自动喷洒控制是提升任务成功率的关键因素，安全控制器能够将碰撞次数降为 0。在整图多区块覆盖任务中，未引导强化学习短训最高仅达到 15.9\% 覆盖率；引入课程式专家引导探索后，DRQN 在 3 个随机种子上达到 100.0\% 成功率和 $97.56\%\pm0.00\%$ 覆盖率，平均路径长度为 $4425.0\pm81.9$ m，碰撞次数和重复喷洒率均为 0。实验表明，复杂农业喷洒任务更适合采用“全局规划先验/课程引导 + DRQN 局部决策”的分层强化学习框架，而非依赖普通 DQN 独立完成全部任务。

\textbf{关键词：} 深度强化学习；DQN；Dueling DQN；DRQN；课程学习；无人机喷洒；路径规划；精准农业；Gazebo 仿真
\end{abstract}

\section{引言}

精准农业要求农业机械能够根据作物分布、地块结构和作业需求进行精细化执行。无人机因具有部署灵活、成本较低、适合中小尺度地块作业等优点，已广泛应用于农药喷洒、巡检、作物长势监测和变量作业等任务。在果园场景中，无人机需要在树行、车辆、人员、建筑物等复杂环境要素之间飞行，并完成对目标树木的有效覆盖。因此，路径规划算法不仅需要保证覆盖率，还需要考虑障碍物安全距离、重复喷洒、路径长度和任务执行效率。

传统喷洒路径通常采用规则行遍历或最近目标贪心等启发式方法。这类方法实现简单，在规则农田或果园环境中具有一定可用性，但当目标分布、障碍物位置和起点条件发生变化时，往往会产生路径冗余和重复喷洒。相比之下，强化学习能够通过与环境交互学习路径规划策略，将覆盖、避障和路径长度等多目标统一纳入奖励函数中，因此适合用于无人机喷洒路径规划问题。

本文提出一种基于深度强化学习的果园无人机喷洒路径规划方法。与仅在抽象网格地图中验证算法不同，本文直接复用 AAS 原有 Gazebo 果园环境，解析仿真世界文件中的目标树木和静态障碍物，并在此基础上构建强化学习训练环境。训练完成后，算法输出的路径可进一步转换为 AAS mission YAML 文件，用于 Gazebo 中的无人机任务执行。

需要说明的是，本文的算法主线并不是将普通 DQN 作为唯一最终算法，而是以 DQN 族价值函数方法为研究基础，逐步比较基础 DQN、Double DQN、Dueling DQN、Rainbow DQN lite 和 DRQN 在不同喷洒任务难度下的表现。普通 DQN 主要用于提供最基础的深度 Q 学习对照；Dueling DQN 在静态树行覆盖任务中表现稳定，可作为静态场景的强基线；DRQN 通过引入循环记忆结构，更适合描述动态障碍、变量需求和长距离转场等部分可观决策过程。因此，本文最终推荐的完整系统主体算法为“DRQN + 分层喷洒控制 + 课程式专家引导探索”，而不是单独的普通 DQN。

本文主要贡献如下：
\begin{itemize}
  \item 构建了一套基于 AAS 原有 \texttt{apple\_orchard.sdf} 的无人机喷洒路径规划实验框架，避免了与仿真系统脱节的手工地图。
  \item 将果园喷洒路径规划建模为强化学习任务，设计了面向目标覆盖、重复喷洒抑制、变量需求满足、静态/动态避障和任务完成效率的状态空间、动作空间和奖励函数。
  \item 对 DQN、Double DQN、Dueling DQN、Rainbow DQN lite、PPO、A2C、SAC、TD3 和 DRQN 进行多随机种子对比实验，并从成功率、覆盖率、需求满足率、路径长度、重复喷洒率、碰撞次数和动态安全值 $S_v$ 等指标进行评价。
  \item 在分层增强任务上补充 5 随机种子主实验、自动喷洒/安全控制/动作空间消融实验和训练学习曲线，用于分析模块贡献和训练稳定性。
  \item 构建了树行、矩形大田和整图多区块三类目标模式，验证同一 AAS 果园地图中从局部树行到大范围农田覆盖的任务扩展能力。
  \item 针对整图稀疏目标导致的探索困难，提出课程式专家引导探索，使 DRQN 在整图多区块任务中达到稳定 97\% 以上覆盖率。
  \item 将路径结果转换为 AAS mission 文件，并生成 Gazebo/QGroundControl 可回放任务和方形动作矩阵图，为三维飞行验证和论文图表展示提供接口。
\end{itemize}

\section{相关工作}

\subsection{无人机农业作业路径规划}

农业无人机路径规划通常需要解决覆盖完整性、飞行安全和作业效率之间的权衡。在规则农田中，行遍历、往复式扫描和最近邻贪心等方法常被用于覆盖规划。这类方法不需要训练，计算成本较低，但对于存在障碍物、目标分布不均匀或起点约束的果园环境，路径往往较长，重复覆盖较多。

图搜索和优化类方法，例如 A*、Dijkstra、人工势场和蚁群算法等，可以在已知地图中搜索可行路径，并通过代价函数考虑路径距离和障碍物风险。然而，这类方法通常依赖显式地图和人工设计的代价函数，在复杂任务约束下调参成本较高。对于喷洒任务而言，路径规划还需要同时考虑覆盖率和重复喷洒率，仅优化距离并不能完全反映喷洒质量。

\subsection{强化学习路径规划}

强化学习通过智能体与环境交互学习策略，适合处理路径规划中的序列决策问题。Q-learning 和 DQN 通过学习状态--动作价值函数实现离散动作选择，已被广泛用于栅格地图导航和避障任务。Double DQN 通过分离动作选择和目标价值估计缓解 Q 值过估计问题；Dueling DQN 将状态价值函数和动作优势函数分离，有助于在多个相似动作之间进行稳定评估；Rainbow DQN 则进一步融合多种 DQN 改进机制。

PPO 属于策略梯度类算法，通过限制策略更新幅度提高训练稳定性。与 DQN 系列不同，PPO 直接优化策略函数，既可用于离散动作任务，也可扩展到连续动作任务。本文将 PPO 与多种 DQN 变体在同一果园喷洒路径规划环境中进行对比，以分析不同强化学习算法在覆盖率、路径长度和训练稳定性方面的差异。

\subsection{本文工作定位}

已有强化学习无人机导航研究多集中于抽象二维环境、室内避障或目标跟踪任务，而较少直接复用真实仿真系统中的农业场景模型。本文重点关注果园喷洒任务中的高层路径规划问题：在 AAS 原有农业场景基础上，将目标树木、农田区块、静态障碍物、变量喷洒需求和动态障碍风险统一抽象为可训练的网格决策环境，并比较不同 DQN 族算法在路径覆盖和喷洒控制中的适用性。

本文不是单纯证明某一个 DQN 算法在所有场景中最优，而是通过递进实验分析算法适用边界。静态树行任务用于验证 DQN 族算法能否完成基本覆盖和避障；智能灌溉任务用于验证变量需求和喷洒控制是否可以纳入奖励函数；动态障碍任务用于验证 $S_v$ 安全值和 DRQN 记忆结构的必要性；整图多区块任务用于验证算法能否从局部树行扩展到更大范围农业地图。这样的定位使本文既有可运行系统实现，也包含失败诊断和改进方案。

从当前实验结果看，普通 DQN 适合作为最基础的价值函数基线，Dueling DQN 适合作为静态覆盖任务的主对照，而 DRQN 更适合作为完整增强系统的主体算法。其原因在于农业喷洒任务并不只是单步避障问题，还包含历史覆盖状态、动态障碍变化、变量需求残差和跨区块转场等长期依赖信息。DRQN 中的循环结构能够利用历史观测缓解部分可观问题；课程式专家引导则用于解决整图稀疏目标下的早期探索困难。因此，本文最终算法主线可概括为“DQN 族对比验证 + DRQN 分层增强 + 课程式全局引导”。

\section{问题建模与方法}

\subsection{环境建模}

本文使用 AAS 工程中的 \texttt{apple\_orchard.sdf} 作为仿真环境来源。程序解析 SDF 中当前启用的 \texttt{include} 模型，并进一步解析嵌套模型文件，例如 \texttt{birch\_row}、\texttt{apple\_grid}、\texttt{apple} 和 \texttt{birch}。当前实验中，\texttt{birch\_row\_1} 为启用状态，\texttt{apple\_grid\_1/2/3} 仍处于注释状态，因此当前目标为已启用树行中的 21 个树木目标。

环境统计如表 \ref{tab:env} 所示。

\begin{table}[H]
\centering
\caption{当前果园仿真环境统计}
\label{tab:env}
\begin{tabular}{lr}
\toprule
项目 & 数值 \\
\midrule
目标树木数量 & 21 \\
静态障碍物数量 & 8 \\
网格分辨率 & 5 m \\
起始网格位置 & $(29, 7)$ \\
当前最佳任务航点数 & 73 \\
\bottomrule
\end{tabular}
\end{table}

规划层将果园环境离散化为网格地图。目标树木所在网格被标记为喷洒目标，障碍物周围一定半径内的网格被标记为不可通行区域。无人机在目标附近喷洒半径内经过时，视为完成对应目标喷洒。

\subsection{马尔可夫决策过程建模}

无人机喷洒路径规划可建模为马尔可夫决策过程：
\begin{equation}
M = (S, A, P, R, \gamma),
\end{equation}
其中 $S$ 为状态空间，$A$ 为动作空间，$P$ 为状态转移函数，$R$ 为奖励函数，$\gamma$ 为折扣因子。

\subsection{状态空间}

为降低训练难度，本文采用 11 维紧凑状态向量，而非直接输入完整地图矩阵：
\begin{equation}
s_t = [r_t, c_t, r_g, c_g, \Delta r, \Delta c, C_t, b_u, b_d, b_l, b_r].
\end{equation}
其中，$r_t,c_t$ 表示无人机当前位置的归一化网格坐标；$r_g,c_g$ 表示最近未覆盖目标的归一化位置；$\Delta r,\Delta c$ 表示无人机与最近目标的相对位置；$C_t$ 表示当前覆盖率；$b_u,b_d,b_l,b_r$ 表示上下左右四个方向是否被障碍物或边界阻挡。

\subsection{动作空间}

对于 DQN、Double DQN、Dueling DQN、Rainbow DQN lite 和 PPO，动作空间定义为四个离散移动动作：
\begin{equation}
A = \{\mathrm{up}, \mathrm{down}, \mathrm{left}, \mathrm{right}\}.
\end{equation}
每次动作使无人机在网格中移动一个单元。若动作导致越界或进入障碍物区域，则位置不更新，并记录一次规划层碰撞惩罚。

\subsection{奖励函数}

奖励函数综合考虑路径长度、目标覆盖、重复喷洒和障碍物风险：
\begin{equation}
R_t = r_{\mathrm{step}} + r_{\mathrm{new}} + r_{\mathrm{repeat}} + r_{\mathrm{progress}} + r_{\mathrm{collision}} + r_{\mathrm{finish}}.
\end{equation}
当前实验中主要奖励项设置如下：
\begin{align}
r_{\mathrm{step}} &= -0.01, \\
r_{\mathrm{new}} &= 5.0 \times N_{\mathrm{new}}, \\
r_{\mathrm{repeat}} &= -0.15 \times N_{\mathrm{repeat}}, \\
r_{\mathrm{collision}} &= -2.0, \\
r_{\mathrm{finish}} &= 50.0.
\end{align}
其中 $N_{\mathrm{new}}$ 表示当前步新增喷洒目标数量，$N_{\mathrm{repeat}}$ 表示重复喷洒目标数量。当覆盖率达到 100\% 时，当前回合结束。

\subsection{主体算法选择与系统框架}

根据任务复杂度，本文将算法角色划分为三个层次。第一层为基础 DQN，用于验证深度 Q 学习是否能够在离散网格中学习“靠近目标、避开障碍、完成覆盖”的基本策略。DQN 的优点是结构简单、易于解释，能够作为所有改进算法的基础对照；但在本文实验中，普通 DQN 对随机种子较敏感，在整图多区块任务中容易陷入局部探索，因此不适合作为最终完整系统的唯一主算法。

第二层为 DQN 改进算法，包括 Double DQN、Dueling DQN 和 Rainbow DQN lite。Double DQN 通过分离动作选择和目标价值估计缓解 Q 值过估计；Dueling DQN 将状态价值 $V(s)$ 与动作优势 $A(s,a)$ 分离：
\begin{equation}
Q(s,a)=V(s)+A(s,a)-\frac{1}{|\mathcal{A}|}\sum_{a'}A(s,a'),
\end{equation}
在多个移动方向价值接近时具有更稳定的动作评估能力。实验表明，Dueling DQN 在静态树行覆盖任务中表现最稳定，因此本文将其作为静态喷洒路径规划的强基线。Rainbow DQN lite 融合了 Double DQN、Dueling 网络和优先经验回放，能在部分任务中提升覆盖效果，但当前实现尚未包含 C51 分布式价值学习和 NoisyNet 探索，因此本文将其定位为增强型对照算法，而不是最终主算法。

第三层为 DRQN（Deep Recurrent Q-Network）及其分层增强版本。DRQN 在 DQN 的价值函数框架中加入 GRU 循环记忆层，使策略能够利用历史观测信息。对于农业喷洒任务，这一点尤其重要：无人机当前局部观测无法完整表达已覆盖区域、动态障碍运动趋势和跨区块转场历史，若仅依靠单步状态，策略容易在局部区域反复移动。DRQN 的循环隐状态可表示为：
\begin{equation}
h_t = \mathrm{GRU}(f(s_t), h_{t-1}), \quad Q_t = g(h_t),
\end{equation}
其中 $f(\cdot)$ 为状态编码网络，$h_t$ 为历史隐状态，$g(\cdot)$ 输出各离散动作的 Q 值。

在完整系统中，本文进一步将 DRQN 与分层喷洒控制和课程式专家引导结合。分层喷洒控制将“是否喷洒”从高维动作空间中部分解耦，通过自动喷洒逻辑减少无效探索；课程式专家引导在训练早期以概率 $p_t$ 采用最短路径专家动作：
\begin{equation}
p_t = p_0 \left(1 - \frac{t}{T_{\mathrm{guide}}}\right),
\end{equation}
当 $t \geq T_{\mathrm{guide}}$ 时 $p_t=0$，策略完全由 DRQN 决策。该设计并不是在测试阶段替代强化学习策略，而是在训练早期提供全局方向，使智能体先学会跨区块转场和目标覆盖，再逐渐依靠自身价值函数完成决策。由此，本文最终推荐的主体算法为 DRQN 分层引导版本；普通 DQN 和 Dueling DQN 分别作为基础基线和静态强基线。

\section{实验设置}

\subsection{对比算法}

本文按任务层级对比不同强化学习算法。静态树行实验主要比较 DQN、Double DQN、Dueling DQN、Rainbow DQN lite 和 PPO；增强实验进一步加入 DRQN；扩展筛选实验补充 A2C、SAC、TD3 和 Maskable PPO 等算法。各算法在本文中的定位如下：
\begin{itemize}
  \item \textbf{DQN：} 基础深度 Q 网络，使用经验回放和目标网络。它是本文的基础对照算法，用于判断最基本的深度 Q 学习是否能够完成覆盖和避障。
  \item \textbf{Double DQN：} 使用在线网络选择动作，目标网络估计目标价值，降低过估计。
  \item \textbf{Dueling DQN：} 将状态价值和动作优势分离，提高相似动作下的价值估计能力。实验中它在静态树行覆盖任务中表现最稳定，因此作为静态任务强基线。
  \item \textbf{Rainbow DQN lite：} 融合 Double DQN、Dueling Network 和 Prioritized Replay 的简化 Rainbow 版本。该版本尚未包含 C51 分布式价值学习和 NoisyNet。
  \item \textbf{PPO：} 基于策略梯度的近端策略优化算法，通过裁剪目标函数限制策略更新幅度。
  \item \textbf{DRQN：} 在 DQN 框架中加入 GRU 循环记忆层，用于处理动态障碍、历史覆盖状态和部分可观条件。本文将其作为完整增强系统和整图多区块任务的主体算法。
  \item \textbf{A2C、SAC、TD3 和 Maskable PPO：} 作为补充强化学习基线，用于判断策略梯度和连续动作算法是否更适合当前离散网格喷洒任务。
\end{itemize}

此外，本文还保留果园行遍历、最近目标贪心和最短路径专家引导作为传统启发式或训练辅助基线。需要强调的是，专家引导只用于训练早期的课程式探索，不在最终评估阶段替代 DRQN 策略。

\subsection{评价指标}

本文使用以下指标评价算法性能：
\begin{itemize}
  \item \textbf{成功率：} 多随机种子实验中达到指定任务阈值且碰撞次数为 0 的比例。静态树行实验以 100\% 覆盖为目标，增强实验以 90\% 或 97\% 需求满足率为目标，整图任务以 97\% 覆盖率为目标。
  \item \textbf{覆盖率：} 已喷洒目标数占全部目标数的比例。
  \item \textbf{需求满足率：} 对智能灌溉任务，已满足剂量需求占总需求的比例。
  \item \textbf{路径长度：} 路径中相邻航点之间欧氏距离的累计值。
  \item \textbf{重复喷洒率：} 重复喷洒次数占总喷洒次数的比例。
  \item \textbf{剂量RMSE：} 变量需求图中目标剂量与实际喷洒剂量之间的均方根误差。
  \item \textbf{最小安全值 $S_v$：} 动态障碍场景中无人机与移动障碍之间的最小归一化安全距离。
  \item \textbf{碰撞次数：} 进入障碍物或非法区域的次数。
  \item \textbf{航点数：} 生成路径中包含的网格航点数量。
\end{itemize}

覆盖率定义为：
\begin{equation}
C = \frac{|T_{\mathrm{sprayed}}|}{|T|} \times 100\%.
\end{equation}
重复喷洒率定义为：
\begin{equation}
R_{\mathrm{repeat}} = \frac{N_{\mathrm{repeat}}}{N_{\mathrm{new}} + N_{\mathrm{repeat}}} \times 100\%.
\end{equation}

\subsection{训练设置}

不同实验采用不同训练预算，以保证实验目的清晰。静态树行任务通常训练 20000 步，并在 3--5 个随机种子上统计均值和标准差；分层增强主实验采用随机种子 7、11、19、23 和 31，每个算法训练 15000 步，最大回合步数为 500，目标阈值为 90\% 需求满足率；消融实验同样采用上述 5 个随机种子，每个配置训练 10000 步；整图多区块优化实验采用 20000 步、3 个随机种子，并开启课程式专家引导探索。训练完成后，对每个模型在同一果园环境中进行确定性评估，并统计覆盖率、需求满足率、路径长度、重复喷洒率、碰撞次数、$S_v$ 和航点数量。实验结果使用均值和标准差表示。

\subsection{统计说明与学习曲线}

本文中多随机种子实验均报告均值和样本标准差。成功率定义为满足任务阈值且规划碰撞次数为 0 的随机种子比例；在智能灌溉任务中，任务进度以需求满足率为主指标，在普通覆盖任务和整图多区块任务中以覆盖率为主指标。除特别说明外，增强任务的最大回合步数为 500，若策略未在最大步数内达到目标，则该回合记为未成功，路径长度通常达到最大步数对应上限。

为观察训练过程而不只比较最终表格结果，本文在 DQN、DQN 变体和 DRQN 训练脚本中增加周期性确定性评估，记录 episode reward、覆盖率和需求满足率随训练步数变化的学习曲线。当前消融实验以每 5000 步评估一次的方式记录 5 个随机种子的均值与标准差；分层增强主实验中，新增的随机种子 23 和 31 已记录学习曲线，若最终论文需要完全一致的 5 随机种子曲线，应使用相同脚本重新强制运行主实验并保存全部曲线。

硬件与软件环境在最终论文定稿时应补充 CPU、GPU、内存、操作系统、Python、PyTorch、Gymnasium、Stable-Baselines3 以及 AAS/Gazebo 版本或提交号，以保证复现实验统计口径完整。

\subsection{分阶段实验设计}

为避免一次性叠加过多任务难点导致实验结果难以解释，本文将实验划分为四个递进层级，如表 \ref{tab:three_stage} 所示。第一级验证基础路径规划能力，第二级单独验证变量需求和喷洒控制，第三级加入动态障碍与 $S_v$ 安全约束，第四级扩展到整图多区块覆盖并引入课程式专家引导，用于验证系统是否能从局部树行扩展到大范围农业地图。

\begin{table}[H]
\centering
\caption{分阶段实验结构}
\label{tab:three_stage}
\resizebox{\textwidth}{!}{
\begin{tabular}{llll}
\toprule
实验层级 & 开启模块 & 主要评价指标 & 实验目的 \\
\midrule
静态果园覆盖实验 & 静态障碍、目标覆盖 & 成功率、覆盖率、路径长度、重复喷洒率、碰撞次数 & 验证基础路径规划与静态避障能力 \\
智能灌溉需求实验 & 变量需求图、喷洒开关 & 需求满足率、剂量RMSE、过喷率、重复喷洒率 & 验证喷洒控制与变量需求匹配能力 \\
完整增强实验 & 动态障碍、$S_v$、变量需求、DRQN & 需求满足率、动态碰撞、最小 $S_v$、路径长度 & 验证动态风险下的完整系统能力与局限 \\
整图多区块覆盖实验 & \texttt{blocks} 目标、自动喷洒、课程引导 & 覆盖率、路径长度、重复喷洒率、碰撞次数 & 验证大范围农田覆盖与长距离转场能力 \\
\bottomrule
\end{tabular}
}
\end{table}

\section{实验结果与分析}

\subsection{多随机种子整体结果}

表 \ref{tab:multiseed} 给出了三随机种子下五种强化学习算法的平均结果。表中 ``$\pm$'' 后为标准差。

\begin{table}[H]
\centering
\caption{五种强化学习算法多随机种子结果}
\label{tab:multiseed}
\resizebox{\textwidth}{!}{
\begin{tabular}{lrrrrr}
\toprule
算法 & 成功率(\%) & 覆盖率(\%) & 路径长度(m) & 重复喷洒率(\%) & 碰撞次数 \\
\midrule
Dueling DQN & 66.7 & $92.1 \pm 13.7$ & $1093.3 \pm 1218.6$ & $62.4 \pm 28.4$ & $0.0 \pm 0.0$ \\
Double DQN & 66.7 & $92.1 \pm 13.7$ & $1110.0 \pm 1203.8$ & $65.9 \pm 27.0$ & $0.0 \pm 0.0$ \\
DQN & 66.7 & $84.1 \pm 27.5$ & $1106.7 \pm 1206.7$ & $62.6 \pm 28.8$ & $0.0 \pm 0.0$ \\
Rainbow DQN lite & 66.7 & $81.0 \pm 33.0$ & $1090.0 \pm 1221.4$ & $69.8 \pm 25.2$ & $0.0 \pm 0.0$ \\
PPO & 33.3 & $54.0 \pm 39.9$ & $1793.3 \pm 1224.0$ & $73.8 \pm 21.9$ & $0.0 \pm 0.0$ \\
\bottomrule
\end{tabular}
}
\end{table}

\begin{figure}[H]
\centering
\includegraphics[width=0.95\textwidth]{outputs/multiseed_20260513_top5/summary/multiseed_summary.png}
\caption{五种强化学习算法在三个随机种子下的均值与标准差对比}
\label{fig:multiseed}
\end{figure}

从表 \ref{tab:multiseed} 和图 \ref{fig:multiseed} 可以看出，Dueling DQN、Double DQN、DQN 和 Rainbow DQN lite 的成功率均为 66.7\%，PPO 成功率为 33.3\%。所有算法在规划层碰撞次数上均为 0，说明当前奖励函数和障碍物网格约束能够有效避免静态障碍物碰撞。

需要注意的是，平均路径长度的标准差较大。这是因为部分随机种子下算法未能达到 100\% 覆盖，回合持续到最大步数限制，导致路径长度达到 2500 m。因此，仅看总体平均路径长度不能完全反映成功策略的路径质量，还需要分析成功回合下的表现。

\subsection{成功回合结果}

表 \ref{tab:successonly} 仅统计达到 100\% 覆盖且碰撞次数为 0 的成功回合。

\begin{table}[H]
\centering
\caption{成功回合中的路径质量对比}
\label{tab:successonly}
\begin{tabular}{lrrrr}
\toprule
算法 & 成功次数 & 路径长度(m) & 重复喷洒率(\%) & 航点数 \\
\midrule
Dueling DQN & 2/3 & $390.0 \pm 42.4$ & $46.8 \pm 12.1$ & $79.0 \pm 8.5$ \\
Double DQN & 2/3 & $415.0 \pm 7.1$ & $51.9 \pm 16.8$ & $84.0 \pm 1.4$ \\
DQN & 2/3 & $410.0 \pm 0.0$ & $46.0 \pm 3.9$ & $83.0 \pm 0.0$ \\
Rainbow DQN lite & 2/3 & $385.0 \pm 35.4$ & $55.6 \pm 7.9$ & $78.0 \pm 7.1$ \\
PPO & 1/3 & $380.0 \pm 0.0$ & $70.4 \pm 0.0$ & $77.0 \pm 0.0$ \\
\bottomrule
\end{tabular}
\end{table}

从成功回合结果看，PPO 在唯一成功回合中路径长度最短，为 380 m，但其成功率仅为 33.3\%，稳定性不足。Rainbow DQN lite 在成功回合中的平均路径长度为 $385.0\pm35.4$ m，路径质量较好，但同样存在训练波动。Dueling DQN 成功率为 66.7\%，成功回合平均路径长度为 $390.0\pm42.4$ m，重复喷洒率为 $46.8\pm12.1\%$，在成功率和路径质量之间取得了较好平衡。

DQN 的成功回合路径长度稳定为 410 m，重复喷洒率为 $46.0\pm3.9\%$，说明基础 DQN 在成功时具有较稳定的路径形态和较低的重复喷洒率，但其平均覆盖率低于 Dueling DQN 和 Double DQN。Double DQN 虽然平均覆盖率与 Dueling DQN 相同，但成功回合路径长度略长，重复喷洒率也略高。

\subsection{单随机种子结果与多随机种子结果差异}

在单次实验中，Rainbow DQN lite 曾取得 100\% 覆盖、0 碰撞和 360 m 路径长度的结果，表现最优。然而，多随机种子实验显示其成功率仅为 66.7\%，且不同随机种子之间波动较大。这说明单次实验结果不能充分代表算法稳定性。相比之下，Dueling DQN 在多个随机种子下表现更均衡，既能保持较高成功率，又能在成功回合中获得较短路径。因此，本文最终更倾向于选择 Dueling DQN 作为当前系统的主算法。

\subsection{智能灌溉需求实验：变量需求与喷洒开关}

在第二级实验中，本文在静态果园环境基础上加入变量喷洒需求图，并开启喷洒开关控制。与基础实验中“经过目标附近即自动喷洒”的设定不同，智能体在该实验中需要同时决定移动方向和是否喷洒，动作空间扩展为 8 个离散动作，即四个移动方向与喷洒开/关的组合。该实验关闭动态障碍，重点考察算法是否能够在不增加动态风险的条件下学习需求匹配策略。

该实验可通过如下配置运行：
\begin{verbatim}
--goal-metric demand --intelligent-irrigation --spray-control
\end{verbatim}
其成功判据为需求满足率达到 97\% 且碰撞次数为 0。相比静态覆盖实验，该实验更关注剂量质量，因此除覆盖率和路径长度外，还需要报告需求满足率、剂量 RMSE、过喷率与重复喷洒率。该层级的意义在于将“喷洒控制是否有效”与“动态避障是否有效”分离，避免在完整增强实验中无法判断性能下降究竟来自变量喷洒还是动态障碍。

\subsection{完整增强实验：动态障碍与智能灌溉}

为进一步贴近原始研究设想，本文在原有果园环境基础上加入了动态障碍、变量喷洒需求图和移动障碍安全值 $S_v$。该增强环境仍然复用 AAS 原有 \texttt{apple\_orchard.sdf}，不额外新建果园地图。与静态覆盖实验相比，增强实验主要增加以下机制：
\begin{itemize}
  \item \textbf{动态障碍：} 在可通行网格中生成 2 个周期性移动障碍，其未来一步位置被用于动作屏蔽和碰撞判定。
  \item \textbf{智能灌溉/变量喷洒需求：} 对目标树木生成 $[0.5,2.0]$ 范围内的需求量，目标由单纯覆盖率扩展为需求满足率。
  \item \textbf{安全值 $S_v$：} 根据无人机与最近动态障碍的曼哈顿距离计算 $S_v\in[0,1]$，数值越大表示越安全。
  \item \textbf{DRQN：} 新增基于 GRU 的循环 DQN，用于在部分可观条件下利用历史观测推断动态障碍趋势。
\end{itemize}

增强实验采用 5 个随机种子（7、11、19、23、31），每个算法训练 30000 步。成功判据定义为需求满足率达到 97\% 且碰撞次数为 0。实验结果见表 \ref{tab:enhanced} 和图 \ref{fig:enhanced}。

\begin{table}[H]
\centering
\caption{动态障碍与智能灌溉增强实验结果（5 个随机种子）}
\label{tab:enhanced}
\resizebox{\textwidth}{!}{
\begin{tabular}{lrrrrrrrr}
\toprule
算法 & 成功率(\%) & 需求满足率(\%) & 覆盖率(\%) & 路径长度(m) & 重复喷洒率(\%) & 动态碰撞 & 最小 $S_v$ & 剂量RMSE \\
\midrule
Rainbow DQN lite & 20.0 & $56.9 \pm 35.4$ & $59.0 \pm 34.9$ & $2094.0 \pm 907.8$ & $57.3 \pm 13.7$ & $0.0 \pm 0.0$ & $1.00 \pm 0.00$ & $1.68 \pm 0.53$ \\
DRQN & 20.0 & $49.4 \pm 47.3$ & $50.5 \pm 48.2$ & $2096.0 \pm 903.4$ & $39.6 \pm 36.3$ & $0.0 \pm 0.0$ & $0.87 \pm 0.30$ & $1.84 \pm 0.68$ \\
Dueling DQN & 0.0 & $40.6 \pm 29.0$ & $43.8 \pm 30.4$ & $2500.0 \pm 0.0$ & $52.4 \pm 34.9$ & $0.0 \pm 0.0$ & $1.00 \pm 0.00$ & $20.77 \pm 43.36$ \\
DQN & 0.0 & $36.5 \pm 19.9$ & $39.0 \pm 21.4$ & $2500.0 \pm 0.0$ & $59.9 \pm 9.7$ & $0.0 \pm 0.0$ & $1.00 \pm 0.00$ & $1.65 \pm 0.62$ \\
\bottomrule
\end{tabular}
}
\end{table}

\begin{figure}[H]
\centering
\includegraphics[width=0.95\textwidth]{outputs/enhanced_dynamic_irrigation_5seeds/summary/multiseed_summary.png}
\caption{动态障碍与智能灌溉增强实验的 5 随机种子统计结果}
\label{fig:enhanced}
\end{figure}

从表 \ref{tab:enhanced} 可以看出，增强任务明显比静态覆盖任务更困难。Rainbow DQN lite 和 DRQN 均仅在 1 个随机种子中达到 97\% 需求满足率，对应成功率为 20.0\%；DQN 和 Dueling DQN 在 30000 步训练下尚未达到成功阈值。所有算法的静态碰撞和动态碰撞均为 0，说明动作屏蔽和 $S_v$ 风险约束能够有效避免碰撞，但覆盖效率和需求满足率仍然不足。

成功回合均出现在随机种子 31。Rainbow DQN lite 在该回合中达到 100.0\% 覆盖率和 97.9\% 需求满足率，路径长度为 470 m，重复喷洒率为 63.2\%；DRQN 同样达到 100.0\% 覆盖率和 97.9\% 需求满足率，路径长度为 480 m，重复喷洒率为 70.0\%。图 \ref{fig:enhanced_path} 给出了成功回合中 Rainbow DQN lite 的路径、需求热力图以及动态障碍轨迹。

\begin{figure}[H]
\centering
\includegraphics[width=0.78\textwidth]{outputs/enhanced_dynamic_irrigation_5seeds/summary/rainbow_seed31_path.png}
\caption{Rainbow DQN lite 在增强任务成功回合中的路径、需求分布与动态障碍轨迹}
\label{fig:enhanced_path}
\end{figure}

上述结果说明，当前代码已经实现了动态障碍、变量喷洒需求、$S_v$ 安全值和 DRQN 训练流程，并且在部分随机种子下能够达到 97\% 需求满足率。但从论文表述角度看，目前不能宣称“稳定达到 97\%”，更准确的结论应为：在增强任务下，Rainbow DQN lite 与 DRQN 已经具备达到 97\% 需求满足率的能力，但 30000 步训练下多随机种子稳定性仍不足，后续需要增加训练步数、改进奖励函数或采用课程学习以提高成功率。

\subsection{优化增强配置尝试与失败诊断}

在前述增强实验基础上，本文进一步尝试了更严格且更接近真实喷洒决策的优化配置：训练步数增加至 100000 步，动态障碍采用果园走廊模式，并加入喷洒开关控制。该配置使动作空间由 4 个移动动作扩展为 8 个“移动 + 喷洒开关”动作，同时动态障碍更接近可行通道，因而对策略探索提出了更高要求。实验仍采用 5 个随机种子和 97\% 需求满足率作为成功阈值。

\begin{table}[H]
\centering
\caption{优化增强配置实验结果（5 个随机种子，100000 步）}
\label{tab:enhanced_optimized}
\resizebox{\textwidth}{!}{
\begin{tabular}{lrrrrrrrr}
\toprule
算法 & 成功率(\%) & 需求满足率(\%) & 覆盖率(\%) & 路径长度(m) & 重复喷洒率(\%) & 动态碰撞 & 最小 $S_v$ & 航点数 \\
\midrule
DRQN & 0.0 & $41.7 \pm 35.4$ & $44.8 \pm 37.1$ & $2500.0 \pm 0.0$ & $47.3 \pm 26.6$ & $0.0 \pm 0.0$ & $0.47 \pm 0.30$ & $501.0 \pm 0.0$ \\
Rainbow DQN lite & 0.0 & $10.1 \pm 14.9$ & $11.4 \pm 15.6$ & $2500.0 \pm 0.0$ & $13.0 \pm 29.2$ & $0.0 \pm 0.0$ & $0.60 \pm 0.37$ & $501.0 \pm 0.0$ \\
\bottomrule
\end{tabular}
}
\end{table}

\begin{figure}[H]
\centering
\includegraphics[width=0.95\textwidth]{outputs/enhanced_optimized_5seeds/summary/multiseed_summary.png}
\caption{优化增强配置下 DRQN 与 Rainbow DQN lite 的 5 随机种子统计结果}
\label{fig:enhanced_optimized}
\end{figure}

表 \ref{tab:enhanced_optimized} 表明，简单增加训练步数并加入喷洒开关并未改善结果，反而使两个算法均未达到 97\% 成功阈值。DRQN 的平均需求满足率为 $41.7\%\pm35.4\%$，明显高于 Rainbow DQN lite，但仍不足以完成完整增强任务；Rainbow DQN lite 在喷洒开关动作空间下出现探索不足，平均需求满足率仅为 $10.1\%\pm14.9\%$。所有回合路径长度均达到 2500 m，说明策略未能在最大步数内完成任务。

该结果说明，完整增强任务的主要瓶颈不是碰撞安全，而是高维动作空间与稀疏任务奖励下的探索效率。后续改进应优先采用课程学习：先在静态覆盖任务中训练路径策略，再在变量需求任务中训练喷洒开关，最后加入动态障碍和 $S_v$ 约束；或者将喷洒开关设计为规则控制器，将强化学习策略集中于路径与避障决策。

\subsection{分层喷洒控制改进实验}

针对优化增强配置中出现的探索困难，本文进一步采用分层控制思想对增强任务进行简化：强化学习策略仅负责上下左右四方向移动，喷洒开关由需求感知规则控制器负责。当当前位置的喷洒半径内存在尚未满足剂量需求的目标单元时，规则层自动开启喷洒；若局部目标需求已经满足，则关闭喷洒。与此同时，安全控制器在动作将导致越界、静态障碍碰撞或动态障碍碰撞时，对动作进行局部安全覆盖。这样可以将智能体从“同时学习移动、避障和喷洒开关”的 8 动作空间中解放出来，使其主要学习路径选择与动态避障。

该实验采用 5 个随机种子（7、11、19、23、31）、每个算法训练 15000 步，并将短跑筛选阈值设置为 90\% 需求满足率。实验开启智能灌溉需求、2 个走廊动态障碍、自动喷洒控制和安全控制器。需要强调的是，该阈值低于前述 97\% 最终目标，其作用是快速判断分层控制是否改善训练方向，而不是替代最终 97\% 实验。实验结果如表 \ref{tab:hierarchical_autospray} 所示。

\begin{table}[H]
\centering
\caption{分层喷洒控制增强实验结果（5 个随机种子，15000 步，90\% 需求目标）}
\label{tab:hierarchical_autospray}
\resizebox{\textwidth}{!}{
\begin{tabular}{lrrrrrrr}
\toprule
算法 & 成功率(\%) & 需求满足率(\%) & 覆盖率(\%) & 路径长度(m) & 重复喷洒率(\%) & 动态碰撞 & 最小 $S_v$ \\
\midrule
DQN & 80.0 & $78.7 \pm 29.3$ & $82.9 \pm 30.4$ & $844.0 \pm 926.5$ & $53.1 \pm 5.4$ & $0.0 \pm 0.0$ & $0.33 \pm 0.00$ \\
DRQN & 60.0 & $80.9 \pm 14.7$ & $83.8 \pm 15.6$ & $1270.0 \pm 1123.2$ & $57.5 \pm 3.5$ & $0.0 \pm 0.0$ & $0.33 \pm 0.00$ \\
Dueling DQN & 60.0 & $55.0 \pm 50.2$ & $57.1 \pm 52.2$ & $1244.0 \pm 1147.1$ & $34.4 \pm 31.5$ & $0.0 \pm 0.0$ & $0.60 \pm 0.37$ \\
Rainbow DQN lite & 20.0 & $60.2 \pm 24.4$ & $63.8 \pm 23.7$ & $2092.0 \pm 912.3$ & $56.7 \pm 4.3$ & $0.0 \pm 0.0$ & $0.33 \pm 0.00$ \\
\bottomrule
\end{tabular}
}
\end{table}

\begin{figure}[H]
\centering
\includegraphics[width=0.95\textwidth]{outputs/hierarchical_autospray_3seeds_15k/summary/multiseed_summary.png}
\caption{分层喷洒控制增强实验的 5 随机种子多算法结果}
\label{fig:hierarchical_autospray}
\end{figure}

\begin{figure}[H]
\centering
\includegraphics[width=0.78\textwidth]{outputs/hierarchical_autospray_3seeds_15k/summary/drqn_seed11_path.png}
\caption{DRQN 在分层喷洒控制增强实验中的代表性路径、需求分布与动态障碍轨迹}
\label{fig:hierarchical_drqn_path}
\end{figure}

从结果可以看出，分层控制相比 8 动作喷洒开关配置明显改善了增强任务的可学习性。在仅训练 15000 步的条件下，DQN 达到 80.0\% 成功率，是当前 5 随机种子短训设置下成功率最高的算法；DRQN 的平均需求满足率最高，为 $80.9\%\pm14.7\%$，且方差低于 DQN，说明循环记忆结构对变量需求和动态障碍具有一定稳定性优势。Dueling DQN 的成功率为 60.0\%，成功回合路径较短，但整体方差较大；Rainbow DQN lite 当前仅达到 20.0\% 成功率，说明轻量 Rainbow 结构在该短训配置下尚不稳定。四种算法的静态碰撞和动态碰撞均为 0，表明自动喷洒与安全控制组合可以在保证安全的同时提高增强任务可学习性。

成功回合单独统计显示，DQN、DRQN、Dueling DQN 和 Rainbow DQN lite 的成功回合平均路径长度分别为 $430.0\pm43.2$ m、$450.0\pm40.0$ m、$406.7\pm50.3$ m 和 $460.0\pm0.0$ m。该结果支持本文的核心改进方向：复杂农业喷洒任务不宜简单扩大动作空间并依赖强化学习一次性学习全部决策，而更适合采用“强化学习路径规划 + 规则喷洒控制 + 局部安全覆盖”的分层框架。后续若继续追求 97\% 稳定成功率，应在该分层框架上增加训练步数、采用课程学习，并进一步降低重复喷洒率。

\subsection{自动喷洒、安全控制与动作空间消融实验}

为进一步分析分层增强任务中各模块的贡献，本文补充分层增强消融实验。所有消融均采用 5 个随机种子（7、11、19、23、31）、10000 训练步、最大回合步数 500、90\% 需求满足率成功阈值，并在 DQN 与 DRQN 之间比较。实验变量包括：是否启用自动喷洒控制、是否启用安全控制器，以及动作空间为 4 个移动动作还是 8 个“移动 + 喷洒开关”动作。表 \ref{tab:ablation_autospray_safety} 给出了每组配置中的最优算法结果。

\begin{table}[H]
\centering
\caption{分层增强任务消融实验结果（5 个随机种子，10000 步）}
\label{tab:ablation_autospray_safety}
\resizebox{\textwidth}{!}{
\begin{tabular}{lllllrrr}
\toprule
配置 & 动作空间 & 喷洒策略 & 安全控制 & 最优算法 & 成功率(\%) & 需求满足率(\%) & 覆盖率(\%) \\
\midrule
Always spray, no safety & 4 动作 & 始终喷洒 & 关闭 & DRQN & 0.0 & $48.2 \pm 31.5$ & $52.4 \pm 33.8$ \\
Auto-spray, no safety & 4 动作 & 自动喷洒 & 关闭 & DRQN & 40.0 & $86.0 \pm 13.7$ & $88.6 \pm 14.9$ \\
Always spray, safety & 4 动作 & 始终喷洒 & 开启 & DQN & 40.0 & $59.8 \pm 40.7$ & $64.8 \pm 42.4$ \\
Auto-spray, safety & 4 动作 & 自动喷洒 & 开启 & DQN & 60.0 & $72.0 \pm 29.4$ & $75.2 \pm 29.8$ \\
Explicit spray, no safety & 8 动作 & 学习开关 & 关闭 & DRQN & 20.0 & $88.1 \pm 9.1$ & $91.4 \pm 11.4$ \\
Explicit spray, safety & 8 动作 & 学习开关 & 开启 & DRQN & 60.0 & $72.1 \pm 27.6$ & $75.2 \pm 27.8$ \\
\bottomrule
\end{tabular}
}
\end{table}

消融结果表明，自动喷洒控制是当前增强任务中最关键的功能模块。在 4 动作设置下，关闭安全控制时由 always-spray 切换为 auto-spray，可使最优成功率从 0.0\% 提高到 40.0\%；开启安全控制后，auto-spray 配置进一步达到 60.0\% 成功率。安全控制器的主要作用是降低风险和消除碰撞：在 auto-spray no-safety 配置中，DRQN 的碰撞和动态碰撞均为 $0.6\pm0.5$，而对应的 auto-spray safety 配置中 DQN 的碰撞和动态碰撞均为 0。

从 DQN 与 DRQN 的对比看，DRQN 在无安全控制或 8 动作显式喷洒开关配置中通常具有更高需求满足率，说明循环记忆有助于处理更复杂的动作--状态耦合；但在 4 动作 auto-spray safety 配置下，DQN 的成功率更高，说明当喷洒控制和局部安全覆盖由规则层承担后，简单价值函数也可以形成较稳定的路径策略。因此，本文将 DQN 作为强基线，将 DRQN 作为更适合复杂部分可观任务和整图扩展的主体候选。

\subsection{训练学习曲线}

除最终评估表外，本文补充训练过程学习曲线，以观察 reward、覆盖率和需求满足率随训练步数变化的趋势。图 \ref{fig:learning_curve_ablation_autospray_safety} 给出了 4 动作 auto-spray safety 消融配置下的 5 随机种子学习曲线，曲线为均值，阴影为标准差。该配置与表 \ref{tab:ablation_autospray_safety} 中成功率最高的 4 动作分层配置对应，因此可用于展示分层控制在训练过程中的收敛趋势。

\begin{figure}[H]
\centering
\includegraphics[width=0.95\textwidth]{outputs/paper_required_experiments/ablation_autospray_safety/summary/learning_curves.png}
\caption{分层增强任务中 auto-spray + safety 配置的训练学习曲线，包括 episode reward、覆盖率和需求满足率。}
\label{fig:learning_curve_ablation_autospray_safety}
\end{figure}

当前消融实验已具备完整 5 随机种子学习曲线；分层增强主实验的最终结果表已补足 5 随机种子，但由于其中 7、11、19 三个随机种子来自较早训练结果，训练时尚未记录学习曲线，因此主实验学习曲线目前仅包含新增随机种子 23 和 31。若最终论文需要严格对应表 \ref{tab:hierarchical_autospray} 的 5 随机种子曲线，应重新以 \texttt{--force --learning-curves} 运行该主实验。

\subsection{扩展强化学习基线筛选}

为进一步判断是否存在比 DQN/DRQN 更适合分层增强任务的强化学习算法，本文在相同分层设置下补充了 A2C、PPO、SAC 和 TD3。实验仍采用 3 个随机种子（7、11、19），训练步数缩短为 10000 步，成功阈值为 90\% 需求满足率。SAC 和 TD3 通过连续动作包装器输出二维方向，再映射为网格移动动作，因此它们作为连续控制基线使用，但与当前离散路径规划任务并非完全同构。

\begin{table}[H]
\centering
\caption{分层增强任务中的扩展强化学习算法筛选结果（3 个随机种子，10000 步）}
\label{tab:extra_rl_screen}
\resizebox{\textwidth}{!}{
\begin{tabular}{lrrrrrr}
\toprule
算法 & 成功率(\%) & 需求满足率(\%) & 覆盖率(\%) & 路径长度(m) & 重复喷洒率(\%) & 动态碰撞 \\
\midrule
DRQN & 66.7 & $75.1 \pm 29.2$ & $77.8 \pm 30.2$ & $1111.7 \pm 1202.8$ & $58.0 \pm 3.9$ & $0.0 \pm 0.0$ \\
DQN & 66.7 & $70.5 \pm 37.4$ & $73.0 \pm 38.5$ & $1116.7 \pm 1198.5$ & $55.8 \pm 5.1$ & $0.0 \pm 0.0$ \\
TD3 & 0.0 & $11.8 \pm 14.5$ & $14.3 \pm 17.2$ & $2500.0 \pm 0.0$ & $13.9 \pm 24.1$ & $0.0 \pm 0.0$ \\
A2C & 0.0 & $11.3 \pm 6.9$ & $14.3 \pm 8.2$ & $2500.0 \pm 0.0$ & $5.6 \pm 9.6$ & $0.0 \pm 0.0$ \\
PPO & 0.0 & $8.5 \pm 1.1$ & $9.5 \pm 0.0$ & $2500.0 \pm 0.0$ & $27.8 \pm 25.5$ & $0.0 \pm 0.0$ \\
SAC & 0.0 & $7.7 \pm 0.7$ & $9.5 \pm 0.0$ & $2500.0 \pm 0.0$ & $11.1 \pm 19.2$ & $0.0 \pm 0.0$ \\
\bottomrule
\end{tabular}
}
\end{table}

\begin{figure}[H]
\centering
\includegraphics[width=0.95\textwidth]{outputs/hierarchical_extra_rl_3seeds_10k/summary/multiseed_summary.png}
\caption{分层增强任务中的扩展强化学习算法筛选结果}
\label{fig:extra_rl_screen}
\end{figure}

结果表明，A2C、PPO、SAC 和 TD3 在短步数筛选中均未达到 90\% 需求满足率。TD3 和 SAC 的安全覆盖次数明显较多，说明连续动作到离散网格动作的映射在该任务中并不理想。相比之下，DRQN 和 DQN 仍保持 66.7\% 成功率，其中 DRQN 在平均需求满足率和成功回合路径长度上略优。因此，在当前分层增强任务中，DRQN 仍是最合适的主算法候选，DQN 可作为结构简单、可解释性较强的主要对照算法。

\subsection{大田区域覆盖扩展}

前述实验默认使用 AAS 场景中当前启用的 \texttt{birch\_row\_1} 作为目标，因此无人机主要围绕果园中一排树木执行喷洒。为进一步贴近大范围农田喷洒任务，本文在不删除原有树行实验的前提下，新增了大田区域覆盖模式。该模式通过 \texttt{--target-mode field} 启用，将 Gazebo 场景中的矩形农田区域离散为规则喷洒目标网格，而不需要额外新建 Gazebo 世界。

当前大田区域设置为：
\begin{verbatim}
--target-mode field --field-bounds=-55,-105,75,-20 --field-spacing 10
\end{verbatim}
其中 \texttt{field-bounds} 表示矩形农田区域的东向和北向边界，\texttt{field-spacing} 表示相邻喷洒需求点间距。该设置在蓝框农田区域内生成 126 个抽象喷洒目标点。基于规则往复式覆盖生成的初始 mission 文件为 \texttt{spray\_field\_lawnmower.yaml}，在规划网格中达到 100.0\% 大田目标覆盖率，路径长度约为 1333.6 m，碰撞次数为 0。

需要说明的是，大田区域扩展与原始树行实验相互独立：默认参数仍为 \texttt{--target-mode trees}，用于复现实验中红框树行目标；只有显式指定 \texttt{--target-mode field} 时，系统才切换到蓝框大田覆盖任务。因此，本文可以同时保留“原始果园树木喷洒”实验和“矩形大田区域覆盖”扩展实验。后续正式强化学习实验可在该 field 模式下继续训练 DRQN 与 DQN，并以农田区域覆盖率、需求满足率、重复喷洒率和路径长度作为主要评价指标。

\subsection{整图多区块覆盖与道路区域处理}

进一步观察 Gazebo 场景可知，蓝框区域并不是新的仿真地图，而是 AAS 原始 \texttt{apple\_orchard.sdf} 世界中的一块矩形农田纹理区域；红框区域对应树行障碍较集中的果园局部。为了避免论文实验只围绕单排树木展开，本文在保留 \texttt{trees} 与 \texttt{field} 两类任务的基础上，新增 \texttt{blocks} 多区块覆盖模式，用若干矩形地块描述整张航拍农田图中的可喷洒区域，并将道路、停车区、建筑物和非农田纹理排除为非喷洒目标。

在 \texttt{blocks} 模式中，评价指标只统计目标农田格子的覆盖情况，因此道路区域不会被计入有效喷洒面积。需要区分的是，QGroundControl 中显示的红色轨迹为无人机飞行轨迹，而不是喷洒沉积图；当多个农田区块彼此不连通时，无人机在区块之间转场时可能从道路或非喷洒区域上方经过。当前代码通过喷洒目标掩膜避免“把道路当作喷洒对象”，但尚未强制禁止飞行轨迹跨越道路。如果后续要求“飞行轨迹也完全不经过道路”，需要将整图任务进一步拆分为多个地块子任务，或在全局规划层加入道路禁飞代价。

为验证整图多区块覆盖接口，本文先构造两个规则基线：最近目标优先覆盖和果园行向往复式覆盖。二者不作为最终强化学习算法，只用于验证 \texttt{blocks} 目标掩膜、道路排除和整图可视化是否正确。结果如表 \ref{tab:whole_farm_rule_baseline} 所示，两个规则基线均可达到 100.0\% 农田目标覆盖率，但路径长度差异较大，说明区块顺序和行向规划会显著影响整图任务效率。

\begin{table}[H]
\centering
\caption{整图多区块覆盖规则基线结果}
\label{tab:whole_farm_rule_baseline}
\begin{tabular}{lrrrrrrr}
\toprule
方法 & 覆盖率(\%) & 路径长度(m) & 上 & 下 & 左 & 右 & 非目标经过率(\%) \\
\midrule
最近目标优先 & 100.0 & 4640.0 & 218 & 213 & 242 & 255 & 90.3 \\
果园行向往复 & 100.0 & 7740.0 & 678 & 704 & 61 & 105 & 92.8 \\
\bottomrule
\end{tabular}
\end{table}

\begin{figure}[H]
\centering
\includegraphics[width=0.88\textwidth]{outputs/action_matrix_square_whole_farm_blocks/action_matrix_square_overlay.png}
\caption{整图多区块覆盖任务的方形动作矩阵可视化。白色为农田目标格，灰色为道路或非喷洒区域，黑色为已覆盖喷洒目标格，橙色为静态障碍，彩色折线表示不同覆盖策略的上、下、左、右动作轨迹。}
\label{fig:whole_farm_blocks_matrix}
\end{figure}

\subsection{方形动作矩阵可视化}

为更直观地展示不同强化学习算法在网格环境中的动作选择，本文新增 \texttt{plot\_action\_matrix.py} 可视化脚本，将算法最终轨迹绘制为接近文献中“方形动作矩阵”的形式。该图不依赖 Gazebo 三维窗口，而是直接读取训练结果中的路径 summary 文件，在二维网格中显示可喷洒目标、已覆盖目标、障碍物、起点以及各算法的离散动作方向。

图 \ref{fig:static_top5_action_matrix} 给出了原始树行静态任务中 5 种算法在同一随机种子下的轨迹矩阵。由于这些模型均在 \texttt{--target-mode trees} 下训练，目标只包含原果园环境中启用的一排树木，因此矩阵中的有效目标区域呈现为较窄的竖向带状区域。这解释了可视化图中“为什么算法总是只跑局部区域”：它不是 Gazebo 地图大小受限，而是当前任务定义只把该树行设为喷洒目标。

\begin{table}[H]
\centering
\caption{静态树行任务中 5 种算法的动作矩阵统计（seed=7）}
\label{tab:static_action_matrix}
\resizebox{\textwidth}{!}{
\begin{tabular}{lrrrrrrr}
\toprule
算法 & 覆盖率(\%) & 路径长度(m) & 上 & 下 & 左 & 右 & 非目标经过率(\%) \\
\midrule
DQN & 100.0 & 82.0 & 25 & 39 & 3 & 15 & 92.8 \\
Double DQN & 100.0 & 360.0 & 40 & 14 & 5 & 13 & 91.8 \\
Dueling DQN & 100.0 & 360.0 & 41 & 15 & 4 & 12 & 90.4 \\
Rainbow DQN lite & 100.0 & 360.0 & 41 & 15 & 4 & 12 & 84.9 \\
PPO & 100.0 & 380.0 & 41 & 15 & 7 & 13 & 76.6 \\
\bottomrule
\end{tabular}
}
\end{table}

\begin{figure}[H]
\centering
\includegraphics[width=0.82\textwidth]{outputs/action_matrix_square_top5_static/action_matrix_square_overlay.png}
\caption{静态树行任务中 DQN、Double DQN、Dueling DQN、Rainbow DQN lite 与 PPO 的方形动作矩阵轨迹。}
\label{fig:static_top5_action_matrix}
\end{figure}

\subsection{整图强化学习短训结果与阶段性失败分析}

在整图多区块任务接口可用后，本文进一步尝试将 DQN、Dueling DQN、Rainbow DQN lite 和 DRQN 直接迁移到 \texttt{blocks} 整图覆盖任务。该组实验采用 1 个随机种子作为快速冒烟验证，训练步数为 8000，最大单回合步数为 2000，目标覆盖阈值为 80\%，并启用自动喷洒控制。命令配置如下：

\begin{verbatim}
--algorithms dqn,dueling-dqn,rainbow-dqn-lite,drqn
--seeds 7 --timesteps 8000 --max-steps 2000
--goal-coverage 0.80 --goal-metric coverage
--target-mode blocks --field-spacing 50 --auto-spray-control
\end{verbatim}

结果如表 \ref{tab:whole_farm_rl_smoke} 所示。该组实验的主要价值不是证明算法已经解决整图喷洒，而是验证当前训练框架能够在更大任务空间下运行，并暴露出直接端到端训练的困难：所有算法均未达到 80\% 覆盖阈值，且路径长度均达到最大步数上限。Rainbow DQN lite 和 Dueling DQN 在短训中略优，但覆盖率仍只有 15.9\% 和 13.4\%；DRQN 和普通 DQN 基本未能有效扩展到整图目标。

\begin{table}[H]
\centering
\caption{整图多区块任务强化学习短训结果（seed=7，8000 步）}
\label{tab:whole_farm_rl_smoke}
\resizebox{\textwidth}{!}{
\begin{tabular}{lrrrrrrrr}
\toprule
算法 & 成功率(\%) & 覆盖率(\%) & 需求满足率(\%) & 路径长度(m) & 上 & 下 & 左 & 右 \\
\midrule
Rainbow DQN lite & 0.0 & 15.9 & 15.9 & 10000.0 & 966 & 958 & 55 & 21 \\
Dueling DQN & 0.0 & 13.4 & 13.4 & 10000.0 & 968 & 963 & 47 & 22 \\
DRQN & 0.0 & 1.2 & 1.2 & 10000.0 & 14 & 0 & 990 & 996 \\
DQN & 0.0 & 0.0 & 0.0 & 10000.0 & 999 & 996 & 0 & 5 \\
\bottomrule
\end{tabular}
}
\end{table}

\begin{figure}[H]
\centering
\includegraphics[width=0.88\textwidth]{outputs/action_matrix_square_whole_farm_rl_smoke_1seed_8k/action_matrix_square_overlay.png}
\caption{整图多区块任务中强化学习算法短训后的动作矩阵轨迹。当前轨迹仍集中在局部区域，说明整图任务不能直接依赖短步数端到端训练解决。}
\label{fig:whole_farm_rl_smoke_matrix}
\end{figure}

该失败结果具有明确的论文价值。首先，它说明前述树行任务中的高覆盖率不能直接外推到整张农田图；当目标从 21 个树行格扩展为多个稀疏区块后，状态空间、动作空间和长距离转场难度显著增加。其次，单纯增加训练步数并不能保证完整增强任务收敛，前文 100000 步增强实验同样没有稳定变好。第三，整图任务更适合采用“全局覆盖规划 + 局部强化学习避障/喷洒控制”的分层方案：规则或图搜索方法先决定区块顺序和主覆盖走廊，DRQN/Dueling DQN 再负责局部避障、动态风险处理、喷洒开关与需求修正。

因此，当前阶段不应将论文主线写成“DQN 完美解决所有农业喷洒任务”，而应表述为：本文在 AAS 果园环境中构建了可运行、可视化、可对比的无人机喷洒路径规划框架；静态树行覆盖任务中 Dueling DQN 表现最稳定；动态障碍与智能灌溉增强任务中 DRQN 是最有潜力的主算法；整图多区块任务已经完成接口、道路排除掩膜和可视化，但直接端到端强化学习仍存在收敛困难，后续需要课程学习、模仿学习或分层规划。

\subsection{课程引导探索后的整图覆盖优化}

针对整图多区块任务中强化学习难以跑满的问题，本文进一步加入课程式专家引导探索。具体做法是在训练早期以一定概率采用最短路径专家动作，将无人机引导至最近的未覆盖目标格；该引导概率随后线性衰减为 0，使后期策略仍由 DRQN 自身价值网络决策。该机制只在显式设置 \texttt{--guided-exploration} 时启用，因此不会改变前述静态树行实验和未引导整图短训实验。

本轮优化采用 DRQN、\texttt{blocks} 整图多区块目标、\texttt{field-spacing=50}、自动喷洒控制、最大步数 1500、目标覆盖率 97\%，训练步数为 20000，并在 3 个随机种子（7、11、19）上验证。结果表明，引导探索显著改善了整图覆盖稳定性：DRQN 在 3 个随机种子上均达到 97.56\% 覆盖率，碰撞次数为 0，重复喷洒率为 0。平均路径长度为 $4425.0\pm81.9$ m，平均航点数为 $886.0\pm16.4$。这说明整图任务的主要瓶颈并非地图接口或目标掩膜错误，而是长距离稀疏目标场景下的探索效率不足。

\begin{table}[H]
\centering
\caption{课程引导探索后 DRQN 在整图多区块任务中的结果（3 个随机种子，20000 步）}
\label{tab:guided_drqn_whole_farm}
\begin{tabular}{lrrrrrr}
\toprule
算法 & 成功率(\%) & 覆盖率(\%) & 路径长度(m) & 重复喷洒率(\%) & 碰撞 & 航点数 \\
\midrule
DRQN + 引导探索 & 100.0 & $97.56\pm0.00$ & $4425.0\pm81.9$ & $0.0\pm0.0$ & $0.0\pm0.0$ & $886.0\pm16.4$ \\
\bottomrule
\end{tabular}
\end{table}

\begin{figure}[H]
\centering
\includegraphics[width=0.88\textwidth]{outputs/action_matrix_square_whole_farm_guided_drqn_3seeds_20k/action_matrix_square_overlay.png}
\caption{课程引导探索后 DRQN 在整图多区块任务中的代表性方形动作矩阵。黑色为已覆盖喷洒目标格，橙色为静态障碍，绿色折线为 DRQN 最终轨迹。}
\label{fig:guided_drqn_whole_farm_matrix}
\end{figure}

与未引导短训相比，该结果将整图覆盖率从最高 15.9\% 提升至 97.56\%，并将最大步数耗尽的局部游走轨迹转化为跨区块覆盖轨迹。该结果进一步支持本文的分层结论：复杂农业喷洒任务需要全局先验或课程引导来解决长距离探索问题，而 DRQN 更适合作为引入历史记忆和局部避障能力的在线策略模块。

\subsection{最终算法选择与论文叙事}

综合上述实验，本文不宜将普通 DQN 简单写成所有任务下的最终最优算法，也不应忽略其在分层增强短训中的强基线价值。普通 DQN 在 5 随机种子 auto-spray + safety 增强任务中取得 80.0\% 成功率，说明当喷洒开关和局部安全覆盖由规则层承担后，基础价值函数方法可以形成有效路径策略；但在显式 8 动作喷洒开关、无安全控制和整图多区块任务中，DQN 的稳定性和探索能力仍然不足。因此，如果论文题目或摘要中直接写“基于 DQN 的完整无人机喷洒系统”，容易被理解为普通 DQN 独立解决所有任务，这与实验结果并不完全一致。

更合理的论文主线是“DQN 族算法对比下的分层喷洒路径规划方法”。在该表述中，DQN 是结构简单且当前增强短训成功率最高的强基线，Dueling DQN 是静态覆盖任务中的稳定基线，DRQN 是面向部分可观、显式喷洒开关和整图扩展任务的主体候选。表 \ref{tab:algorithm_role} 总结了各算法在本文中的角色。

\begin{table}[H]
\centering
\caption{本文中不同强化学习算法的角色定位}
\label{tab:algorithm_role}
\resizebox{\textwidth}{!}{
\begin{tabular}{llll}
\toprule
算法 & 实验表现 & 论文角色 & 是否作为最终主体 \\
\midrule
DQN & 静态任务可用，分层增强短训成功率最高，但整图任务覆盖率低 & 基础深度 Q 学习强基线，用于证明问题可由价值函数方法建模 & 部分是 \\
Double DQN & 静态任务与 DQN 接近，缓解过估计 & DQN 改进对照，说明估计方式变化的影响 & 否 \\
Dueling DQN & 静态树行覆盖最稳定 & 静态果园覆盖任务强基线，可作为第一阶段主结果 & 部分是 \\
Rainbow DQN lite & 部分实验较好，但稳定性不足 & 增强型 DQN 对照，展示优先经验回放等机制的影响 & 否 \\
PPO/A2C/SAC/TD3 & 当前离散网格任务中优势不明显 & 补充基线，说明并非所有 RL 算法都适合该任务 & 否 \\
DRQN & 平均需求满足率较高，在显式喷洒开关和课程引导整图任务中表现更稳 & 完整系统主体候选，负责历史信息利用、局部决策和复杂任务扩展 & 是 \\
\bottomrule
\end{tabular}
}
\end{table}

因此，本文建议的算法表述为：以 DQN 族强化学习为基础，先通过普通 DQN、Double DQN 和 Dueling DQN 验证静态果园覆盖路径规划的可行性；再通过分层自动喷洒和安全控制降低增强任务探索难度，使 DQN 和 DRQN 均能形成有效策略；最后引入 DRQN、课程式专家引导和全局覆盖先验，实现整图多区块农业场景下的稳定覆盖。这样的叙事既保留了 DQN 论文方向，又能真实反映实验结果，避免把普通 DQN 过度拔高或把单一算法写成所有配置下的绝对最优。

如果论文题目需要更准确，可采用以下表述之一：
\begin{itemize}
  \item \textbf{基于DQN族深度强化学习的果园无人机喷洒路径规划方法研究}
  \item \textbf{基于分层DRQN的果园无人机喷洒覆盖路径规划方法研究}
  \item \textbf{面向精准农业喷洒任务的DQN族强化学习路径规划与避障方法研究}
\end{itemize}
其中第一个题目最稳妥，既保留 DQN 方向，又允许正文中把 DRQN 作为最终增强算法；第二个题目更聚焦最终方法，但需要在论文中弱化普通 DQN 的标题地位；第三个题目适合突出算法对比和系统框架。

\subsection{当前代码状态与复现实验指令}

截至当前阶段，代码已经保留并支持三类任务：原始树行喷洒任务 \texttt{trees}、矩形大田覆盖任务 \texttt{field}、整图多区块覆盖任务 \texttt{blocks}。其中 \texttt{trees} 用于复现红框果园树行实验，\texttt{field} 用于蓝框矩形农田演示，\texttt{blocks} 用于接近整张地图的多地块覆盖实验。系统还实现了变量需求图、喷洒开关、动态障碍、移动障碍安全值 $S_v$、DRQN、分层自动喷洒控制、课程式专家引导探索、多算法对比和方形动作矩阵可视化。

常用复现实验命令如下。静态树行多算法对比：
\begin{verbatim}
wsl --cd /mnt/c/Users/zyy/Documents/GitHub/aerial-autonomy-stack bash -lc ".venv/bin/python experiments/spray_dqn/run_top5_multiseed.py --algorithms dqn,double-dqn,dueling-dqn,rainbow-dqn-lite,ppo --seeds 7,11,19,23,31 --timesteps 20000 --output-dir experiments/spray_dqn/outputs/multiseed_5seeds_top5"
\end{verbatim}

分层增强任务 5 随机种子主实验：
\begin{verbatim}
wsl --cd /mnt/c/Users/zyy/Documents/GitHub/aerial-autonomy-stack bash -lc ".venv/bin/python experiments/spray_dqn/run_top5_multiseed.py --algorithms dqn,drqn,dueling-dqn,rainbow-dqn-lite --seeds 7,11,19,23,31 --timesteps 15000 --max-steps 500 --goal-coverage 0.90 --goal-metric demand --dynamic-obstacles 2 --dynamic-obstacle-mode corridor --intelligent-irrigation --auto-spray-control --safety-controller --learning-curves --eval-freq 5000 --output-dir experiments/spray_dqn/outputs/hierarchical_autospray_3seeds_15k"
\end{verbatim}

分层增强任务消融实验：
\begin{verbatim}
wsl --cd /mnt/c/Users/zyy/Documents/GitHub/aerial-autonomy-stack bash -lc ".venv/bin/python experiments/spray_dqn/run_paper_experiments.py --skip-main --output-root experiments/spray_dqn/outputs/paper_required_experiments --ablation-timesteps 10000 --eval-freq 5000"
\end{verbatim}

整图多区块任务推荐从较短的验证命令开始：
\begin{verbatim}
wsl --cd /mnt/c/Users/zyy/Documents/GitHub/aerial-autonomy-stack bash -lc ".venv/bin/python experiments/spray_dqn/run_top5_multiseed.py --algorithms dqn,dueling-dqn,rainbow-dqn-lite,drqn --seeds 7 --timesteps 8000 --max-steps 2000 --goal-coverage 0.80 --goal-metric coverage --target-mode blocks --field-spacing 50 --auto-spray-control --output-dir experiments/spray_dqn/outputs/whole_farm_blocks_rl_smoke_1seed_8k --force"
\end{verbatim}

整图多区块 DRQN 课程引导优化实验：
\begin{verbatim}
wsl --cd /mnt/c/Users/zyy/Documents/GitHub/aerial-autonomy-stack bash -lc ".venv/bin/python experiments/spray_dqn/run_top5_multiseed.py --algorithms drqn --seeds 7,11,19 --timesteps 20000 --max-steps 1500 --goal-coverage 0.97 --goal-metric coverage --target-mode blocks --field-spacing 50 --auto-spray-control --guided-exploration 0.85 --guided-exploration-fraction 0.75 --output-dir experiments/spray_dqn/outputs/whole_farm_blocks_guided_drqn_3seeds_20k --force"
\end{verbatim}

若要生成论文中的方形动作矩阵图，可运行：
\begin{verbatim}
wsl --cd /mnt/c/Users/zyy/Documents/GitHub/aerial-autonomy-stack bash -lc ".venv/bin/python experiments/spray_dqn/plot_action_matrix.py --input-dir experiments/spray_dqn/outputs/whole_farm_blocks_guided_drqn_3seeds_20k --output-dir experiments/spray_dqn/outputs/action_matrix_square_whole_farm_guided_drqn_3seeds_20k --algorithms drqn --style square-overlay --max-arrows 360"
\end{verbatim}

若要由成功 RL 路径生成 Gazebo replay mission，可运行：
\begin{verbatim}
wsl --cd /mnt/c/Users/zyy/Documents/GitHub/aerial-autonomy-stack bash -lc ".venv/bin/python experiments/spray_dqn/prepare_gazebo_replay.py --summary experiments/spray_dqn/outputs/hierarchical_autospray_3seeds_15k/seed_7/metrics/drqn_summary.json --mission-out aircraft/aircraft_resources/missions/spray_paper_replay.yaml --report-out experiments/spray_dqn/outputs/gazebo_replay/gazebo_replay_report.md"
\end{verbatim}

\section{Gazebo 任务生成}

训练完成后，算法输出的网格路径可转换为 AAS mission YAML 文件。当前系统已生成多组可在 AAS/Gazebo 中回放的任务文件：
\begin{itemize}
  \item \texttt{aircraft/aircraft\_resources/missions/spray\_dqn.yaml}
  \item \texttt{aircraft/aircraft\_resources/missions/spray\_rainbow\_dqn.yaml}
  \item \texttt{aircraft/aircraft\_resources/missions/spray\_drqn\_hierarchical.yaml}
  \item \texttt{aircraft/aircraft\_resources/missions/spray\_drqn\_hierarchical\_red\_row.yaml}
  \item \texttt{aircraft/aircraft\_resources/missions/spray\_field\_lawnmower.yaml}
  \item \texttt{aircraft/aircraft\_resources/missions/spray\_whole\_farm\_blocks.yaml}
  \item \texttt{aircraft/aircraft\_resources/missions/spray\_paper\_replay.yaml}
\end{itemize}

对于最终选定的算法，也可以通过结果 summary 文件重新生成 mission。生成的 mission 文件包含起飞、等待、速度设置、航点 reposition 和降落步骤，可由 AAS mission node 在 Gazebo 中执行。实际三维飞行画面可在 Gazebo Sim 窗口中观察，轨迹和遥测状态可在 QGroundControl 中查看。其中 \texttt{spray\_drqn\_hierarchical\_red\_row.yaml} 对应原始红框树行区域，\texttt{spray\_field\_lawnmower.yaml} 对应蓝框矩形农田区域，\texttt{spray\_whole\_farm\_blocks.yaml} 对应整图多区块覆盖演示，\texttt{spray\_paper\_replay.yaml} 对应由分层增强 DRQN 成功路径转换得到的论文回放案例。

表 \ref{tab:gazebo_replay_case} 给出了当前已转换为 Gazebo mission 的 RL 路径案例。该案例来自分层自动喷洒与安全控制增强任务中的 DRQN，随机种子为 7，在规划层达到 95.2\% 覆盖率和 90.9\% 需求满足率，满足 90\% 需求阈值且无规划碰撞和动态碰撞。

\begin{table}[H]
\centering
\caption{由 RL 路径生成的 Gazebo replay mission 案例}
\label{tab:gazebo_replay_case}
\begin{tabular}{ll}
\toprule
指标 & 数值 \\
\midrule
算法与随机种子 & DRQN, seed 7 \\
Mission 文件 & \texttt{spray\_paper\_replay.yaml} \\
覆盖率 & 95.2\% \\
需求满足率 & 90.9\% \\
路径长度 & 410.0 m \\
规划碰撞 & 0 \\
动态碰撞 & 0 \\
最小 $S_v$ & 0.33 \\
Mission 航点数 & 83 \\
\bottomrule
\end{tabular}
\end{table}

当前工作已经完成从 RL 规划路径到 AAS mission YAML 的转换，并生成了 Gazebo replay 命令与验证清单。正式论文提交前，还应补充 Gazebo/QGroundControl 实际执行证据，包括仿真世界启动截图、无人机起飞与航点执行日志、ROS 2 位置话题记录、任务完成或降落日志，以及无可见碰撞的画面证据。这样可以将“规划层可行路径”进一步闭环为“三维仿真可执行任务”。

\section{结论与展望}

本文基于 AAS 原有 Gazebo 果园仿真环境，构建了无人机喷洒路径规划强化学习实验框架，并对 DQN、Double DQN、Dueling DQN、Rainbow DQN lite、PPO 和 DRQN 等算法进行了分阶段对比。静态覆盖实验表明，Dueling DQN 在基础树行任务中表现稳定，适合作为静态覆盖强基线。进一步的动态障碍与智能灌溉增强实验表明，显式 8 动作喷洒开关会显著增加探索难度，即使延长训练步数也难以稳定收敛；采用自动喷洒控制和安全控制器后，分层增强任务在 5 个随机种子上明显改善，其中 DQN 达到 80.0\% 成功率，DRQN 达到最高平均需求满足率 $80.9\%\pm14.7\%$，四种 DQN 族算法的规划碰撞和动态碰撞均为 0。消融实验进一步证明，自动喷洒控制是提升需求满足率和成功率的关键模块，安全控制器则主要负责消除碰撞风险。本文还补充了大田矩形区域、整图多区块覆盖、道路排除掩膜、Gazebo mission 回放和方形动作矩阵可视化，使实验从“单排树行覆盖”扩展为可讨论整张农业地图覆盖能力的系统框架。在整图多区块任务中，未引导强化学习短训最高仅达到 15.9\% 覆盖率；加入课程式专家引导探索后，DRQN 在 3 个随机种子上达到 100.0\% 成功率和 $97.56\%\pm0.00\%$ 覆盖率，说明该任务需要全局先验或课程引导来突破长距离探索瓶颈。

当前研究仍存在一定局限。首先，动态障碍和智能灌溉目前实现于 Python 高层规划环境中，Gazebo 内尚未生成真实移动行人或车辆，也尚未模拟真实药液雾化、风场漂移和沉积分布。其次，分层增强实验目前采用 90\% 需求满足率作为阶段性成功阈值，虽然已能证明 auto-spray 与 safety-controller 的有效性，但若要支撑更强的“高精度需求满足”结论，仍需补充 97\% 阈值下的更长训练或课程学习实验。再次，整图多区块任务虽然在课程引导后取得了稳定覆盖结果，但该结果仍依赖最短路径专家在训练早期提供示范，因此更适合被解释为“全局规划先验 + DRQN 局部策略”的分层框架，而不是纯端到端 DQN 独立解决整图喷洒。最后，当前已经生成由成功 RL 路径转换得到的 \texttt{spray\_paper\_replay.yaml} mission，但正式论文提交前仍应补充实际 Gazebo/QGroundControl 回放截图与日志证据。后续应采用更系统的课程学习、模仿学习、全局覆盖规划与局部 DRQN/Dueling DQN 避障控制相结合的分层策略，以提升动态增强任务中的多随机种子稳定性和真实 Gazebo 场景泛化能力。

\end{document}
